\section{Síntesis y ampliación}

Esta mesa redonda se desarrolló justo después del lanzamiento de O1; los tres invitados, desde distintas dimensiones (la perspectiva emprendedora, la perspectiva académica y la perspectiva técnica), emitieron un juicio sobre el estado actual y el futuro de la AGI. El consenso central es el siguiente:

\begin{enumerate}
\item \textbf{Aceleración, no desaceleración}: en los últimos 18 meses el desarrollo de la AGI se ha acelerado, con una doble mejora en cantidad y calidad
\item \textbf{O1 es un cambio de paradigma}: abre una nueva dimensión de RL Scaling, dota al modelo de capacidad de System 2 y impulsa a la AGI de L1 a L2
\item \textbf{La generalización es el reto central}: pasar de Narrow a General requiere resolver los problemas de la definición de reward y de los datos de supervisión de proceso
\item \textbf{La demanda de cómputo estallará}: el lado de inferencia crecerá con certeza, el lado de entrenamiento crecerá con alta probabilidad y el número de parámetros podría reiniciar el Scaling
\item \textbf{Las oportunidades de emprendimiento se amplían}: hay una dirección definida (RL Scaling) pero una ruta incierta, lo que favorece a los equipos innovadores
\item \textbf{Ruta de tres etapas}: simular el mundo, explorar el mundo, inducir el mundo; en los últimos 18 meses ha habido avances emblemáticos en cada etapa
\end{enumerate}

\begin{knowledgebox}{Lecturas adicionales}
\begin{itemize}
\item Blog técnico de OpenAI O1: «Learning to Reason with LLMs», trabajo pionero del test-time scaling
\item Definición de niveles de AGI de OpenAI (L1--L5)
\item Análisis técnico de Tesla FSD V12: modelo grande de conducción autónoma de extremo a extremo
\item AlphaGo $\rightarrow$ AlphaFold $\rightarrow$ AlphaGeometry: la ruta de evolución de RL de DeepMind
\item Kahneman, «Thinking, Fast and Slow»: la obra original de la teoría de System 1/System 2
\end{itemize}
\end{knowledgebox}

%% --- Fin del contenido --- %%

\end{document}
